\documentclass[border=4pt]{standalone}
\usepackage{amsmath,amssymb,mathtools,xcolor,varwidth}
\definecolor{auxpurple}{HTML}{8E24AA}
\begin{document}
\begin{varwidth}{28em}
\large\bfseries
Here is the device. While $B$ approaches $A$, we always \emph{produce} $AB$ and $AD$ onward to      fixed remote points $b$ and $d$, and parallel to the secant $BD$ we draw \textcolor{auxpurple}{the line $bd$}.      Upon the distant ends we erect \textcolor{auxpurple}{the auxiliary arc $Acb$}, always kept \emph{similar} to      the shrinking arc $ACB$. This is the trick: the original figure dwindles toward $A$, but the      enlarged copy out at $b,d$ stays a fixed, finite size --- a faithful magnification in which every      proportion of the little figure is preserved exactly. So whatever ratio arc, chord, and tangent      bear to one another below, the same ratio is held, full-sized and visible, in \textcolor{auxpurple}{the      finite triangle $Abd$ and arc $Acb$} above.
\end{varwidth}
\end{document}
